\documentstyle[% 


righttag% 


,11pt% 


,amssymb% 


,russian%               remove in Eng. 


,izvru%                 Set izven% in Eng. 


]{amsart} 


 


% {=}


\newcommand{\re}{\operatorname{Re}} 


\newcommand{\tts}[1]{}  


\renewcommand{\baselinestretch}{0.97}


 


\theoremstyle{definition} 


\theorembodyfont{\rm} 


\newtheorem{notenonum}{\hskip\parindent Замечание} 


\renewcommand{\thenotenonum}{} 


 


 


%\hoffset=-15mm%                  For Eng. 


\setcounter{page}{3} 


\god{2003} 


\nomer{10~(497)} %                        Remove in Eng. 


%\nomer{Vol.\,47, No.\,10}%               Set in Eng. 


%\str{\pageref{begin}--\pageref{end}}%  Set in Eng. 


\udk{517.958} 


 


\author{\it В.М.\,Брук, В.А.\,Крысько} 


% NB:   ^^ remove in Eng. 


\title{О существовании и единственности решения  


задачи Коши для дифференциально-операторных 


уравнений смешанного типа} 


\thanks{} 


 


\date{} 


%\pagestyle{myheadings}%         Set in Eng. 


\pagestyle{plain} %              Remove in Eng. 


\begin{document} 


%\thispagestyle{plain}%          Set in Eng. 


\maketitle 


\label{begin} 


 


 


B работе устанавливаются достаточные условия существования и 


единственности 


решения задачи Коши для системы дифференциально-операторных 


уравнений смешанного типа (уравнения гиперболического и параболического 


типов). При соответствующем подборе операторных коэффициентов из 


рассматриваемого уравнения получается дифференциальное уравнение, 


описывающее колебания пластинки, --- модифицированное уравнение 


Жермен--Лагранжа (гиперболическое уравнение). При этом для определения 


температурного поля используется трехмерное уравнение теплопроводности 


(параболическое уравнение). Полученная система описывает связанную 


задачу термоупругости пластинок. 


 


Пусть $H$ --- гильбертово пространство, $L$, $M$ --- самосопряженные 


отрицательно определенные операторы в $H$ с областями определения 


$D(M)\subset D(L)$, 


$C$ --- ограниченный оператор в $H$. Предполагается, что оператор $L$ 


перестановочен с $C$ и резольвентой $M$ (определение перестановочности 


см., напр., в [1], с.\,322). На отрезке $[0,T]$ рассмотрим систему 


\begin{equation} 


\begin{cases} 


W''(t)+L^2W(t)+\alpha LC\theta(t)=q_1(t),\\ 


\theta'(t)-M\theta(t)-\beta C^*LW'(t)=q_2(t) 


\end{cases} 


\end{equation} 


($\alpha,\beta>0$) с неизвестными функциями $W$, $\theta$. 


 


Выберем, например, в качестве $H=L_2(\Omega)$ пространство функций, 


определенных 


в ограниченной с гладкой границей области $\Omega= \Omega_1\times 


[-h/2,h/2]\subset R^3$ 


($\Omega_1\subset R^2$, $h>0$) и имеющих суммируемый квадрат 


нормы; в качестве $M$, $L$ --- операторы, порожденные соответственно 


выражением $\lambda^2_0\frac{\partial^2}{\partial x^2}+ \lambda^2_0 


\frac{\partial^2}{\partial y^2}+ \frac{\partial^2}{\partial z^2}$ 


($\lambda_0\ne0$) 


с граничным условием $\theta|_{\partial\Omega}=0$ и 


выражением $\gamma(\frac{\partial^2}{\partial x^2}+ 


\frac{\partial^2}{\partial y^2})$ 


($\gamma>0$) с граничным условием $W|_{\partial\Omega_1}=0$. Оператор 


$C$ определим формулой $C\theta=\int\limits^{h/2}_{-h/2} 


z\theta(x,y,z)dz$. Тогда $C^*g=z\int\limits^{h/2}_{-h/2} g(x,y,z)dz$. 


При таком подборе $L$, $M$, $C$ и соответствующем выборе чисел 


$\alpha$, $\beta$, $\gamma$ из (1) получаются известные уравнения 


колебания 


пластинки ([2], с.\,336) и 


уравнение теплопроводности ([3], с.\,12). 


 


Далее в целях сокращения записи считаем, что в (1) $\alpha=\beta=1$. В 


противном 


случае рассуждения практически не изменятся. 


 


Для свед\'ения системы (1) к одному уравнению первого порядка введем 


действующий в пространстве $\wt H=H\times H\times H$ оператор 


$$ 


A_0= 


\begin{pmatrix} 


0 & E & 0 \\ -L^2 & 0 & -LC \\ 0 & C^*L & M 


\end{pmatrix} 


. 


$$ 


Здесь и далее $E$ обозначает тождественный оператор в соответствующем 


пространстве. Отметим, что перестановочность операторов $C$ и $L$ влечет 


перестановочность $C^*$ и $L^*=L$ ([1], с.\,323). 


 


 


\begin{lem} 


Оператор $A_0$ допускает замыкание $A$ в пространстве $\wt H$, а 


оператор $A^{-1}$ существует, ограничен и определен на всем пространстве 


$\wt H$. 


\end{lem} 


 


 


\begin{pf} 


Операторная матрица 


$$ 


\begin{pmatrix} 


CM^{-1}C^* & -L^{-2} & -L^{-1}CM^{-1}\\ 


E & 0 & 0 \\ 


-LM^{-1}C^* & 0 & M^{-1} 


\end{pmatrix} 


$$ 


определяет ограниченный оператор в $\wt H$. Непосредственная проверка 


показывает, что ее произведение (в любом порядке) 


на операторную матрицу, определяющую оператор $A_0$, 


дает единичную операторную матрицу. Отсюда следуют 


все утверждения леммы, причем оператор $A^{-1}$ определяется последней 


операторной матрицей. 


\end{pf} 


 


 


В дальнейшем оператор и его замыкание будут обозначаться одним символом. 


B случае незамкнутости оператор заменяется своим замыканием. 


 


Если в (1) положить $v=W'$ и через $u$ обозначить столбец 


$(W,v,\theta)^\top$, то в 


пространстве $\wt H$ однородную систему (1) можно записать в виде 


\begin{equation} 


u'(t)=Au(t). 


\end{equation} 


 


Напомним ([4], сс.\,38, 76), что решением уравнения (2) 


называется функция, удовлетворяющая (2) на замкнутом отрезке $[0,T]$, а 


ослабленным решением (2) называется функция, непрерывная на $[0,T]$ и 


удовлетворяющая (2) на $(0,T]$, причем на $(0,T]$ ее производная 


предполагается 


непрерывной. 


 


 


Обозначим через $\wt L{}^{-2}$ оператор в $\wt H$, определяемый равенством 


$$ 


\wt L{}^{-2}= 


\begin{pmatrix} 


L^{-2} & 0 & 0 \\ 0 & L^{-2} & 0 \\ 0 & 0 & L^{-2} 


\end{pmatrix} 


. 


$$ 


Из свойств перестановочных операторов следует, что $\wt L{}^{-2}$ и


$A^{-1}$ перестановочны. 


 


 


\begin{thm} 


Задача Коши для уравнения $(2)$ с начальным условием 


\begin{equation} 


u(0)=u_0 


\end{equation} 


имеет единственное решение, если $u_0\in A^{-2}\wt L{}^{-2}(\wt H)$, 


и единственное ослабленное решение, если $u_0\in A^{-1}\wt L{}^{-4}(\wt H)$. 


\end{thm} 


 


 


{\bf Доказательство} основано на приводимых ниже оценках резольвенты 


оператора $A$. Далее решение строится по 


резольвенте так же, как в ([4], гл.\,1, \S\,3, 


с.\,78), с помощью обратного преобразования Лапласа. 


 


Пусть $\lambda=\sigma+i\tau$. Определим операторы 


$\Delta_1=\Delta_1(\lambda)$, $\Delta_2=\Delta_2(\lambda)$ в $H$ 


формулами 


\begin{align*} 


\Delta_1&=(\lambda^2E+L^2)-\lambda L^2C 


(M-\lambda E)^{-1}C^*,\\ 


\Delta_2&=(M-\lambda E)-\lambda L^2 


(L^2+\lambda^2E)^{-1}C^*C. 


\end{align*} 


 


 


\begin{lem} 


При $\sigma=\re\lambda>0$ операторы $\Delta_1$, $\Delta_2$ имеют 


ограниченные обратные и


$$


\|\Delta^{-1}_1\|\le \frac{1}{\sigma|\lambda|}, \quad


\|\Delta^{-1}_2\|\le\frac{1}{\sigma}. 


$$


\end{lem} 


 


 


\begin{pf} 


Обозначим $S=\frac{1}{\lambda}\Delta_1$. Тогда для любого $f\in D(L^2)$ 


имеем 


\begin{multline*} 


(Sf,f)=(\sigma+i\tau)(f,f)+\frac{1}{|\lambda|^2} 


(\sigma-i\tau)(L^2f,f)- \\ 


-((M-\sigma E+i\tau E)|M-\lambda E|^{-1}LC^*f, 


\ |M-\lambda E|^{-1}LC^*f). 


\end{multline*} 


Отсюда и из отрицательной определенности $M$ вытекает 


$\re(Sf,f)\ge\sigma\|f\|^2$. 


Следовательно, $\|Sf\|\ge\sigma\|f\|$. Аналогично получим 


$\|S^*g\|\ge\sigma\|g\|$. Поэтому оператор $S^{-1}$ 


существует, ограничен, определен на всем $H$ и 


$\|S^{-1}\|\le\frac{1}{\sigma}$. Отсюда следует 


утверждение леммы об операторе $\Delta_1$. 


 


Из равенства $-\lambda L^2(L^2+\lambda^2E)^{-1}= -L^2(\lambda 


L^2+|\lambda|^2\ov\lambda E)|L^2+\lambda^2E|^{-2}$ вытекает 


$\re(-\lambda L^2(L^2+\lambda^2E)^{-1} C^*Cg,g)\le0$ ($g\in H$). 


Отсюда и из отрицательной определенности оператора $M$ 


получим $\re(\Delta_2f,f)\le-\sigma(f,f)$ ($f\in D(\Delta_2)$). 


Поэтому $|\re(\Delta_2f,f)|\ge \sigma\|f\|^2$ и, следовательно, 


$\|\Delta_2f\|\ge\sigma\|f\|$. Аналогично 


получим $\|\Delta^*_2f\|\ge\sigma\|f\|$. Два последних неравенства влекут 


утверждение 


леммы об операторе $\Delta_2$. 


\end{pf} 


 


 


\begin{lem} 


При $\re\lambda>0$ оператор $A$ имеет резольвенту


$$ 


R(\lambda)= 


\begin{pmatrix} 


-\lambda\Delta^{-1}_1+L^2\Delta^{-1}_1C(M-\lambda E)^{-1}C^* & 


-\Delta^{-1}_1 & -\Delta^{-1}_1LC(M-\lambda E)^{-1} \\ 


% 


L^2\Delta^{-1}_1 & -\lambda\Delta^{-1}_1 & 


-\lambda\Delta^{-1}_1LC(M-\lambda E)^{-1} \\ 


% 


-L^3(L^2+\lambda^2E)^{-1}\Delta^{-1}_2C^* & 


\lambda L(L^2+\lambda^2E)^{-1}\Delta^{-1}_2C^* & \Delta^{-1}_2 


\end{pmatrix} 


. 


$$ 


\end{lem} 


 


 


\begin{pf} 


Непосредственной проверкой с учетом перестановочности $L$ и $M^{-1}$ 


можно убедиться, что произведение $R(\lambda)(A_0-\lambda E)$ дает 


единичную операторную матрицу. Кроме того, область значений оператора 


$A_0-\lambda E$ плотна в $\wt H$. Действительно, пусть найдется такой 


элемент 


$x=(x_1,x_2,x_3)\in\wt H$, что $(A_0-\lambda E)^*x=0$. Это равенство 


равносильно следующим трем равенствам: 


$$ 


-\ov\lambda x_1-L^2x_2=0,\quad 


x_1-\ov\lambda x_2+CLx_3=0,\quad 


-C^*Lx_2+(M-\ov\lambda E)x_3=0. 


$$ 


Выражая из первого равенства $x_1$, из последнего $x_3$ и подставляя во 


второе, 


получим $\Delta_1(\ov\lambda)x_2{=}0$. Из леммы 2 вытекает, что $x_2=0$. 


Следовательно, $x_1=x_3=0$. 


Поэтому $x=0$. 


\end{pf} 


 


 


\begin{lem} 


При достаточно больших $\sigma=\re\lambda>0$ справедливы оценки 


\begin{equation} 


\|R(\lambda)\|\le k|\lambda|/\sigma,\quad


\|R(\lambda)\wt L{}^{-2}\|\le k/|\lambda|.


\end{equation} 


Здесь и далее $k$ обозначает постоянную $($не зависящую от $\lambda)$,


различную, вообще говоря, в различных неравенствах. 


\end{lem} 


 


 


\begin{pf} 


Установим сначала неравенство 


\begin{equation} 


\|L^2\Delta^{-1}_1\|\le|\lambda|/\sigma. 


\end{equation} 


Так же, как в доказательстве леммы 2, имеем при $g\in H$ 


\begin{multline*} 


(SL^{-2}g,g)=(\sigma+i\tau)(L^{-2}g,g)+\frac{1}{|\lambda|^2} 


(\sigma-i\tau)(g,g)- \\ 


% 


-((M-\sigma E-i\tau E)|M-\lambda E|^{-1} 


C^*g,\ |M-\lambda E|^{-1}C^*g). 


\end{multline*} 


Следовательно, $\re(SL^{-2}g,g)\ge\sigma/|\lambda|^2$. Отсюда получаем 


$\|L^2S^{-1}\|\le|\lambda|^2/\sigma$; поэтому (5) выполняется. 


 


Обозначим через $a_{ij}$ элементы операторной матрицы, определяющей 


$R(\lambda)$. Тогда 


\begin{alignat*}{3} 


\|a_{11}\|&\le \frac{k}{\sigma}, &\quad 


\|a_{12}\| &\le \frac{1}{|\lambda|\sigma}, &\quad 


\|a_{13}\|&\le \frac{k}{|\lambda|\sigma}, \\ 


% 


\|a_{21}\|&\le \frac{|\lambda|}{\sigma}, &\quad 


\|a_{22}\|&\le \frac{1}{\sigma}, &\quad 


\|a_{23}\|&\le \frac{k}{\sigma},\\ 


% 


\|a_{11}L^{-2}\|&\le \frac{k}{|\lambda|}, &\quad 


\|a_{12}L^{-2}\|& \le \frac{k}{|\lambda|\sigma}, &\quad 


\|a_{13}L^{-2}\| &\le \frac{k}{|\lambda|^2\sigma}, \\ 


% 


\|a_{21}L^{-2}\|&\le \frac{1}{|\lambda|\sigma}, &\quad 


\|a_{22}L^{-2}\|&\le \frac{k}{|\lambda|}, &\quad 


\|a_{23}L^{-2}\|&\le \frac{k}{|\lambda|\sigma}. 


\end{alignat*} 


Эти оценки элементов первых двух строк матрицы $R(\lambda)$ сразу же 


получаются 


из (5), леммы 2, неравенств $\|L(M-\lambda E)^{-1}\|\le k$, 


\begin{equation} 


\|(M-\lambda E)^{-1}\|\le\frac{1}{|\lambda|}, 


\end{equation} 


а также из следующих двух равенств: 


$$ 


a_{11}=-\frac{1}{\lambda}E-\frac{1}{\lambda}L^2 


\Delta^{-1}_1,\quad a_{22}L^{-2}=-\frac{1}{\lambda}L^{-2} 


+\frac{1}{\lambda}\Delta^{-1}_1+\Delta^{-1}_1 


C(M-\lambda E)^{-1}C^*, 


$$ 


которые доказываются элементарными преобразованиями. 


Для оценки элементов третьей строки матрицы $R(\lambda)$ установим 


некоторые 


вспомогательные неравенства. 


Так как $\re\big((\frac{1}{\lambda}L^2+\lambda E)f,f\big)\ge \sigma(f,f)$ 


для любого $f\in D(L^2)$, то $\big\|\big(\frac{1}{\lambda} L^2+\lambda 


E\big)^{-1}\big\|\le\frac{1}{\sigma}$. 


Отсюда вытекает 


\begin{equation} 


\|(L^2+\lambda^2E)^{-1}\|\le\frac{1}{|\lambda|\sigma}. 


\end{equation} 


Из этого неравенства элементарными преобразованиями получаем 


\begin{equation} 


\|L^2(L^2+\lambda^2E)^{-1}\|\le|\lambda|/\sigma. 


\end{equation} 


 


Докажем, что 


\begin{equation} 


\|L\Delta^{-1}_2\|\le m, 


\end{equation} 


где $m$ не зависит от $\lambda$. Для любого $f\in D(\Delta_2L^{-1})$ имеем 


$$ 


(\Delta_2L^{-1}f,f)=(ML^{-1}f,f)-(\sigma+i\tau) 


(L^{-1}f,f)-(\sigma-i\tau)|\lambda|^2(Lg,g)- 


(\sigma+i\tau)(L^3g,g), 


$$ 


где $g=\sqrt{|L^2+\lambda^2E|^{-2}C^*Cf}$. 


Из отрицательной определенности $L$ следует 


$\re(\Delta_2L^{-1}f,f)\ge \re(ML^{-1}f,f)\ge k\|f\|^2$. Последнее 


неравенство вытекает из положительной 


определенности ограниченного оператора $LM^{-1}$, обратного к 


$ML^{-1}$. Таким образом, (9) выполняется. 


 


Равенство $\Delta^{-1}_2L^{-1}=(M-\lambda E)^{-1}L^{-1}+ 


(L\Delta^{-1}_2)(L^2+\lambda^2E)^{-1} C^*C\lambda(M-\lambda E)^{-1}$ 


и неравенства (6), (7), (9) влекут 


\begin{equation} 


\|\Delta^{-1}_2L^{-1}\|\le k/|\lambda|. 


\end{equation} 


Теперь неравенства (7)--(10) и лемма 2 позволяют получить 


оценки элементов третьей строки матрицы $R(\lambda)$ 


\begin{alignat*}{3} 


\|a_{31}\|&\le k\frac{|\lambda|}{\sigma}, &\quad 


\|a_{32}\|&\le \frac{k}{\sigma}, &\quad 


\|a_{33}\|&\le \frac{1}{\sigma}, \\ 


% 


\|a_{31}L^{-2}\| &\le \frac{k}{|\lambda|\sigma}, &\quad 


\|a_{32}L^{-2}\|&\le \frac{k}{|\lambda|\sigma}, &\quad 


\|a_{33}L^{-2}\|&\le \frac{k}{|\lambda|}. 


\end{alignat*} 


Из полученных оценок $a_{ij}$ следует утверждение леммы. 


\end{pf} 


 


 


Отметим, что некоторые оценки можно улучшить. Например, 


$\|a_{12}L^{-2}\|\le\frac{k}{|\lambda|^2}$, однако для дальнейшего этого не 


требуется. 


 


Переходим непосредственно к доказательству теоремы. Следуя ([4], 


с.\,78), решение задачи Коши (2), (3) ищем в виде 


$$ 


u(t)= 


-\frac{1}{2\pi i}\int^{\sigma+i\infty}_{\sigma-i\infty}e^{\lambda t} 


R(\lambda)u(0)d\lambda, 


$$ 


где интегрирование ведется по прямой, параллельной мнимой оси и 


проходящей через точку $(\sigma,0)$ ($\sigma>0$ достаточно большое). 


 


Пусть $u_0\in A^{-2}\wt L{}^{-2}(\wt H)$, т.\,е. $u_0=R^2(0)\wt 


L{}^{-2}x_0$ ($x_0\in \wt H$, $R(0)=A^{-1}$). Из резольвентного тождества 


имеем 


\begin{multline} 


u(t)= 


-\frac{1}{2\pi i}\int^{\sigma+i\infty}_{\sigma-i\infty}e^{\lambda t} 


R(\lambda)R^2(0)\wt L{}^{-2}x_0d\lambda= \\ =


-\frac{1}{2\pi i}\int^{\sigma+i\infty}_{\sigma-i\infty}e^{\lambda t} 


\bigg(\frac{1}{\lambda^2}R(\lambda)-\frac{1}{\lambda^2}R(0)- 


\frac{1}{\lambda}R^2(0)\bigg)\wt L{}^{-2}x_0d\lambda. 


\end{multline} 


Существование интеграла в (11) вытекает из леммы 4. Из той же леммы и из 


равенства 


$\frac{1}{2\pi i}\int^{\sigma+i\infty}_{\sigma-i\infty} \frac{1}{\lambda} 


e^{\lambda t}d\lambda=1$ 


следует непрерывность $u(t)$ на $[0,T]$ и равенство 


$u(0)=R^2(0)\wt L{}^{-2}x_0$. Кроме того, $u(t)$ имеет непрерывную на 


$[0,T]$ производную 


$$ 


u'(t)= 


-\frac{1}{2\pi i}\int^{\sigma+i\infty}_{\sigma-i\infty}e^{\lambda t} 


\bigg(\frac{1}{\lambda}R(\lambda)\wt L{}^{-2}x_0 


-\frac{1}{\lambda}R(0)\wt L{}^{-2}x_0\bigg)d\lambda 


$$ 


и удовлетворяет уравнению (2). Таким образом, функция (11) является 


решением задачи Коши (2), (3). 


 


Пусть теперь $u_0\in A^{-1}\wt L{}^{-4}(\wt H)$, т.\,е. $u_0=R(0)\wt 


L{}^{-4}z_0$ ($z_0\in \wt H$). Рассмотрим функцию 


\begin{equation} 


\wt u(t)= 


-\frac{1}{2\pi i}\int^{\sigma+i\infty}_{\sigma-i\infty}e^{\lambda t} 


R(\lambda)R(0)\wt L{}^{-4}z_0d\lambda= 


-\frac{1}{2\pi i}\int^{\sigma+i\infty}_{\sigma-i\infty}e^{\lambda t} 


\bigg(\frac{1}{\lambda}R(\lambda)\wt L{}^{-4}z_0- 


\frac{1}{\lambda}R(0)\wt L{}^{-4}z_0\bigg)d\lambda. 


\end{equation} 


Из леммы 4 вытекает существование интеграла, непрерывность $\wt u(t)$ на 


$[0,T]$ и 


равенство $\wt u(0)=R(0)\wt L{}^{-4}z_0$. Применяя к первому интегралу в 


(12) формулу 


интегрирования по частям, с учетом леммы 4 получим 


$$ 


\wt u(t)= 


\frac{1}{2\pi i}\frac{1}{t} 


\int^{\sigma+i\infty}_{\sigma-i\infty}e^{\lambda t} 


R^2(\lambda)R(0)\wt L{}^{-4}z_0d\lambda. 


$$ 


Отсюда следует, что $\wt u(t)$ имеет на $(0,T]$ непрерывную производную 


$$ 


\wt u'(t)= 


-\frac{1}{2\pi i}\frac{1}{t^2} 


\int^{\sigma+i\infty}_{\sigma-i\infty}e^{\lambda t} 


R^2(\lambda)R(0)\wt L{}^{-4}z_0d\lambda+ 


\frac{1}{2\pi i}\frac{1}{t} 


\int^{\sigma+i\infty}_{\sigma-i\infty}e^{\lambda t} 


\lambda R^2(\lambda)R(0)\wt L{}^{-4}z_0d\lambda 


$$ 


и удовлетворяет уравнению (2). 


 


Единственность решения вытекает из оценки (4) и из результатов, 


изложенных в ([4], с.\,81). Отметим также, что из той же оценки и из 


([4], с.\,49) следует существование решения задачи (2), (3) при $u_0\in 


D(A^4)$. $\quad\square$


 


 


\begin{notenonum} 


Для функций (11), (12) выполняются неравенства 


$\|u(t)\|\le k\|R(0)x_0\|$, $\|\wt u(t)\|\le k\|\wt L{}^{-2}z_0\|$. Это 


вытекает из леммы 4 и резольвентного 


тождества, примененного в первом интеграле (11). 


\end{notenonum} 


 


 


Перейдем к рассмотрению неоднородного уравнения 


\begin{equation} 


u'(t)=Au(t)+f(t). 


\end{equation} 


Обозначим через $U(t)$ оператор 


$$ 


U(t)x= 


-\frac{1}{2\pi i}\int^{\sigma+i\infty}_{\sigma-i\infty}e^{\lambda t} 


R(\lambda)x\,d\lambda,\quad 0\le t\le T, 


$$ 


определенный при всех $x\in\wt H$, для которых последний интеграл 


существует. Из 


леммы 4 и равенства 


$$ 


U(t)R(0)\wt L{}^{-2}z= 


-\frac{1}{2\pi i}\int^{\sigma+i\infty}_{\sigma-i\infty}e^{\lambda t} 


\frac{1}{\lambda}(R(\lambda)-R(0))\wt L{}^{-2}z\,d\lambda 


\quad (z\in\wt H) 


$$ 


следует, что $U(t)R(0)\wt L{}^{-2}$ является при фиксированном $t$ 


ограниченным оператором 


в $\wt H$, а операторная функция $U(t)R(0)\wt L{}^{-2}$ сильно непрерывна 


на $[0,T]$. 


Из доказательства теоремы 1 вытекает, что при $u_0\in R^2(0)\wt 


L{}^{-2}(\wt H)$ ($u_0\in R(0)\wt L{}^{-4}(\wt H)$) 


функция $U(t)u_0$ является решением (ослабленным решением) задачи Коши 


(2), (3). 


 


 


\begin{thm} 


Пусть функция $f(t)$ удовлетворяет одному из следующих 


условий\/${:}$ 


 


$1)$ $f(t)\in R^2(0)\wt L{}^{-2}(\wt H)$ при любом $t$ и функция $\wt 


L{}^2A^2f(t)$ непрерывна$;$ 


 


$2)$ $f(t)$ имеет производную такую, что $f'(t)\in R(0)\wt L{}^{-2}(\wt 


H)$ при любом $t$, 


функция $\wt L{}^2Af'(t)$ непрерывна, и $f(0)\in R(0)\wt L{}^{-2}(\wt H)$. 


 


Тогда функция 


\begin{equation} 


y(t)=\int^t_0 U(t-s)f(s)ds 


\end{equation} 


является решением уравнения $(13)$. 


\end{thm} 


 


 


\begin{pf} 


Пусть $f$ удовлетворяет условию 1). Обозначим 


$g(t)=\wt L{}^2A^2f(t)$. Тогда функция (14) и ее производная будут иметь 


вид 


\begin{gather*} 


y(t)=\int^t_0 U(t-s)R^2(0)\wt L{}^{-2}g(s)ds, \\ 


% 


y'(t)=R^2(0)\wt L{}^{-2}g(t)+\int^t_0 


U'_t(t-s)R^2(0)\wt L{}^{-2}g(s)ds= \\ 


=f(t)+\int^t_0 AU(t-s)R^2(0)\wt L{}^{-2}g(s)ds=f(t)+Ay(t). 


\end{gather*} 


В силу непрерывности $g$ все интегралы в этих равенствах существуют. 


Следовательно, $y(t)$ --- решение (13). 


 


Пусть теперь $f$ удовлетворяет условию 2). Тогда из равенства 


$f(t)=f(0)+\int\limits^t_0 f'(s)ds$ и свойств $U(t)$ вытекает 


существование интеграла в (14). Несложными 


преобразованиями с использованием формулы интегрирования по 


частям так же, как в ([4], с.\,164), получим 


$$ 


y(t)=\int^t_0 U(t-s)f(s)ds=\int^t_0 U(t-s)R(0) 


f'(s)ds-R(0)f(t)+U(t)R(0)f(0). 


$$ 


Отсюда, вычисляя $y'(t)$, так же, как в первой части доказательства, 


получим, 


что $y(t)$ --- решение (13). 


\end{pf} 


 


 


 


\begin{thebibliography}{9} 


 


\bibitem{1} 


Рисс Ф., С\"екефальви-Надь Б. {\em Лекции по функциональному анализу}. -- 


2-е изд. -- М.: Мир, 1979. -- 588 с. 


\bibitem{2} 


Бабаков И.М. {\em Теория колебаний}. -- М.: Наука, 1965. -- 559 с. 


\bibitem{3} 


Беляев Н.М., Рядно А.А. {\em Методы теории теплопроводности}. -- М.: 


Высш. школа, 1982. -- 325~с. 


\bibitem{4} 


Крейн С.Г. {\em Линейные дифференциальные уравнения в банаховом 


пространстве}. -- М.: Наука, 1967. -- 464 с. 


 


\end{thebibliography} 


 


\vskip 1cm 


 


\postup{\it Саратовский государственный}{\qquad \it Поступили} 


\postup{\it технический университет}{\it первый вариант $28.03.2000$} 


\postup{}{\it окончательный вариант $22.05.2003$} 


 


\label{end} 


\end{document} 


